\chapter{Discussion}
\label{ch:discussion}

This thesis compares supervised classification (Track~1), representation learning (Track~2), and downstream probing of learned representations (Track~3).
While the same backbone architecture is used throughout for consistency, the tracks answer different questions and should be interpreted accordingly.

\paragraph{Linear probe vs.\ end-to-end classification}
Track~3 is implemented as a linear evaluation protocol: an encoder pretrained in Track~2 is reused and (by default) frozen, and only a lightweight classifier head is trained to predict the Chagas label.
This probes whether label information is readily accessible in the learned representation without allowing the encoder to co-adapt to the downstream objective.
In contrast, Track~1 trains encoder and classifier jointly and therefore reflects the best achievable supervised classifier under the chosen objective and preprocessing.
Consequently, Track~1 and Track~3 are not an apples-to-apples comparison of classifier performance; Track~3 is primarily used to connect improvements in embedding-space diagnostics to downstream separability.

\paragraph{Clinical label noise and short-window ECGs.}
The labels in this thesis are based on seropositivity, yet clinical staging defines the indeterminate form (Stage A) as seropositive with a normal ECG and early cardiomyopathy (Stage B1) as the first appearance of ECG abnormalities in a seropositive patient \parencite{chagas_progression_2022,lancet_clinical_chagas_features_2024}.
This creates biological label noise for ECG-based classification because a large fraction of positives may be ECG-normal, while the model's signal is driven by electrophysiological manifestations such as RBBB/LAFB and ventricular ectopy that are typical of CCC \parencite{chagas_ecg_abnormalities_2018,lancet_clinical_chagas_features_2024}.
In addition, the largest dataset (CODE-15\%) uses self-reported Chagas labels, which introduces extra uncertainty beyond staging effects, whereas SaMi-Trop provides clinically confirmed positives \parencite{ribeiro_code-15_2021,ribeiro_sami-trop_2021}.
Conversely, PTB-XL provides comparatively clean negatives because it was recorded in Germany (1989--1996), a non-endemic setting, so Chagas labels are treated as negatives on epidemiological grounds \parencite{wagner_ptb_xl}.
Given short (~6.25~s) windows, the models are inherently biased toward beat morphology and high-burden ectopy rather than low-frequency paroxysmal events, which helps explain why performance is stronger in cohorts enriched for cardiomyopathy and why ranking-based screening is the most clinically plausible objective.

\paragraph{Augmentation strength as a limiting factor.}
The augmentation suite used in this thesis was intentionally mild (small gain jitter and low-variance Gaussian noise), with stronger perturbations such as time warping, time masking, and channel masking disabled to preserve clinical interpretability.
While this choice stabilizes supervised training and supports direct waveform interpretation, it may be suboptimal for contrastive representation learning: Track~2 objectives generally benefit from stronger or more diverse view generation to encourage invariances.
Therefore, some of the observed limits in Track~2 (and downstream Track~3) performance could reflect underpowered augmentations rather than intrinsic limitations of the objectives.

\paragraph{Lead construction and rotation fidelity.}
All experiments operate on the 12-lead signals as provided by the source datasets; no new leads are re-derived from raw electrode potentials, and the only signal modifications are the defined preprocessing and augmentations.
Consequently, the axis-rotation augmentation relies on a fixed linear lead model applied to potentially heterogeneous acquisition pipelines and gain calibrations.
This means the rotation parameters are approximate and may not correspond to a precise physical change in cardiac orientation for every record, which is another reason the rotation should be interpreted as a controlled robustness probe rather than a clinically exact transformation.

\paragraph{Lead selection in ST-DFT-LRP.}
Beat-level ST-DFT-LRP aggregation in this thesis is restricted to Lead~II, which is a standard rhythm lead and provides a clear P--QRS--T morphology.
This choice improves feasibility and interpretability but is not optimal for all Chagas-related abnormalities: several structural and conduction manifestations (e.g., bundle-branch patterns or regional ST--T changes) are known to be more prominent in precordial or alternative limb leads.
Therefore, the lead-II focus should be interpreted as a rhythm-centric view of model attention rather than a comprehensive 12-lead characterization of Chagas cardiomyopathy.

\paragraph{RBBB annotations across datasets.}
For CODE15, RBBB flags are provided directly in the dataset metadata; for PTB-XL, the closest available analogue used here is \texttt{ptb\_crbbb} derived from \texttt{scp\_codes}.
This mapping allows coarse, cross-dataset comparisons of bundle-branch related patterns, but should be interpreted cautiously because complete RBBB in PTB-XL is not a perfect substitute for the broader RBBB labels in CODE15.

\paragraph{Architecture-specificity and generalization.}
All results in this thesis are conditional on a fixed ResNet-18 backbone, chosen to isolate the effects of preprocessing and objective functions rather than to optimize architecture.
This implies that the reported gains and failure modes are most directly applicable to CNN-style inductive biases that emphasize local morphology and translation invariance.
Architectures with different temporal aggregation (e.g., attention-based models or dilated TCNs) may emphasize rhythm, interval structure, or long-range dependencies differently, which can change the sensitivity to augmentation, normalization, and dataset bias. \todo{List 1--2 alternative architecture families relevant to ECGs used in your literature scan.}
In particular, the axis-rotation augmentation may interact differently with models that fuse leads late or model inter-lead correlations explicitly. \todo{Note any lead-fusion strategy differences observed in prior ECG architectures.}
Future work should therefore validate the key findings (e.g., loss-function effects, preprocessing regimes, rotation robustness) on at least one non-CNN baseline to assess architectural transferability. \todo{Specify a concrete baseline to test (e.g., 1D Transformer, InceptionTime, or a dilated TCN).}
Future work should revisit augmentation strength in Track~2 specifically, exploring moderate temporal perturbations or structured masking while validating that ECG morphology remains interpretable and physiologically plausible.

\todo{beobachtungen nur für diese resnet18 architektur?}
